\section{A changing-policy case in a shared knowledge workspace}\label{sec:C24:policycase}
Consider an enterprise research agent preparing an internal recommendation about a vendor. It retrieves a pricing document, a legal assessment, and an operational incident note. The documents have different owners and access scopes. During the task, the legal assessment is corrected and the requesting employee changes teams, losing access to part of the workspace. The agent's active context still contains the earlier excerpts. A fresh retrieval permission check alone cannot remove those excerpts from the current computation or from previously generated summaries.

The first design requirement is dependency tracking. Every extracted claim and generated summary should identify the source versions it used. When the legal assessment changes, the system marks dependent claims stale and prevents them from being promoted as current evidence. This does not imply that all prior artifacts must be automatically deleted: some may be retained as historical records under appropriate access controls. The key distinction is between historical provenance and current applicability.

The second requirement is authority revalidation. A pending action to share the recommendation must be checked under the employee's current principal and purpose. Historical access does not authorize a new disclosure. If the current context contains content no longer permitted for the intended action, the application can rebuild a restricted context or require a controlled review path. It should not ask the model to infer access from the prose, because the decisive policy state may be absent or deliberately untrusted.

The third requirement is correction propagation. Updating the original document and vector index does not update a separately stored summary, a procedural note, or another agent's copied context. A dependency graph can enumerate these derivatives, but actual propagation may be asynchronous. The service should expose whether invalidation is pending, complete, or impossible for a particular external copy. A claim of immediate deletion across all surfaces would require evidence about those surfaces, not simply successful deletion of one database row.

A useful experiment stages the policy change at several points: before retrieval, after retrieval but before synthesis, after synthesis but before sharing, and after a recipient has accessed the artifact. The expected control differs at each stage. Before sharing, the system can block or regenerate. After disclosure, it may revoke further access but cannot guarantee that the recipient forgot or did not copy the information. Recovery metrics should reflect this irreversibility instead of treating every revocation as full restoration.

The comparison conditions should include connector-only checks, retrieval-time checks, and dependency-aware checks spanning derived memory and artifact publication. Hold the model and documents fixed. Measure unauthorized disclosure, use of stale claims, valid recommendation quality, review burden, and completion latency. A system that blocks every recommendation after any policy change has strong containment but poor utility; a system that always completes may have the opposite profile. Reporting both axes prevents a degenerate policy from appearing universally superior.

This case also clarifies why enterprise deployment evidence cannot be borrowed from a generic question-answering benchmark. The difficult behavior is not recalling a sentence but maintaining a coherent relationship among evidence version, principal, purpose, and action time. A memory engine can improve retrieval while leaving that relationship unenforced. A workflow platform can enforce approvals while presenting stale evidence to reviewers. A model can reason correctly about the visible documents while lacking the current policy state. The system must join these layers explicitly.

For implementation, represent each artifact with a content digest, source-version dependencies, audience, purpose, and approval record. Store access decisions separately from the artifact text so they can expire or be revoked. Generate the model-visible authority summary from current policy, and record the epoch used at each effect boundary. These are concrete engineering requirements that can be checked without claiming that the model has a complete understanding of organizational intent. The remaining semantic judgment concerns whether the recommendation follows from the authorized, current evidence; that judgment still needs a calibrated evaluator and appropriate review.

The evaluator should also inspect who could observe each intermediate artifact. A final recommendation may be properly restricted while a temporary spreadsheet or trace exposes the underlying assessment to a broader audience. Thus the privacy unit is the complete artifact lineage, not merely the final answer. Instrumented tests can enumerate created objects and compare their access policies with the task contract. Where an application does not expose intermediate-object visibility, the corresponding claim remains unverified. This is a concrete observation limit and should affect the deployment assurance case rather than be replaced by a general statement that the platform is secure.
